\documentclass[10pt,a4paper]{amsart}
\usepackage[margin=19mm]{geometry}
\usepackage{amsmath,amssymb,mathtools}
\usepackage[scr=rsfs,cal=euler]{mathalfa}
\usepackage{xfrac}
\usepackage[unicode,hidelinks,bookmarks=false]{hyperref}
\newcommand{\ie}{i.e.\ }
\newcommand{\cf}{cf.\ }
\newcommand{\resp}{respectively\ }
\newcommand{\Subsection}{\subsection}
\providecommand{\nobreakdash}{\nobreak\hbox{-}}
\setlength{\emergencystretch}{2em}
\multlinegap0pt
\usepackage{tikz}
\usetikzlibrary{cd}
\usepackage{breqn}
\usepackage{amssymb}
\usepackage{tikz}
\usepackage[shortlabels]{enumitem}
\usepackage[normalem]{ulem}
\usepackage{xcolor}
\newtheorem{lemma}{Lemma}
\newtheorem{theorem}{Theorem}
\newcommand{\F}{\mathcal{F}}
\newcommand{\Hom}{\operatorname{Hom}}
\title{Globalization of Partial Group Actions on Not Necessarily Associative Algebras and Covariant Representations}
\author{Mikhailo Dokuchaev, Emmanuel Jerez,, Jos\'e L. Vilca-Rodr\'iguez}
\date{Source: arXiv:2604.21167v1 (CC BY 4.0)}
\begin{document}
\maketitle

\noindent\emph{Source and licence.} Globalization of Partial Group Actions on Not Necessarily Associative Algebras and Covariant Representations. Version arXiv:2604.21167v1 (CC BY 4.0), distributed by arXiv under the Creative Commons Attribution 4.0 International licence, http://creativecommons.org/licenses/by/4.0/, as recorded in the arXiv licence metadata for this exact version and shown on the abstract page https://arxiv.org/abs/2604.21167v1. Full licence notice: \texttt{licenses/arxiv-2604.21167-CC-BY-4.0.txt}.\par\smallskip

\noindent\textit{Editorial scope.} Counted unit: the article's theorem theorem (source0.tex lines 1236--1245). The article's complete original proof of it is delivered with it (source0.tex lines 1247--1343, one \texttt{proof} environment of 97 lines). No mathematics is written, repaired or re-derived: the statement and the proof are the article's own lines, in the article's own notation, and no retained line is substituted or re-flowed.  Retained uncounted as the source-local context the proof needs: lines 1120--1122; lines 1141--1146.  Nothing was omitted from any retained span, so the package carries no omission record.  No retained line contains a citation, so the wrapper carries no bibliography and no bibliography entry.  The retained lines contain the article's own diagram code, which the wrapper typesets with the standard tikz packages; no image file is delivered and the package declares no asset dependency.  Editorial formatting only, in addition: an AMS \texttt{amsart} preamble that restates the article's own declarations and macros used by the retained lines, namely the environments \texttt{lemma}, \texttt{theorem} on separate counters, together with the standard packages those lines need.  The retained environments print in this document as Lemma 1, Theorem 1 in order of appearance, whereas the article prints the same items with the numbering of its own full document.  The complete native source of the article as distributed in the arXiv e-print for this exact version is bundled unchanged as \texttt{sources/199001/source0.tex}.
\begin{equation}\label{eq xw}
    xw = T(xw) = xT(w) \quad \text{for all } x \in A \text{ and } w \in W.
\end{equation}

\begin{lemma}\label{l: T for Lambda M}
    Let \(\lambda: G \to \operatorname{End}_0(A)\) be  an ideal partial representation  of \( G \) on \(A\) and \((\pi, M)\) be a covariant representation of \(\lambda\). Then the mapping \( T_\Lambda: \Lambda(M) \to \Lambda(M) \) given by \( \lfloor h, m \rfloor \mapsto \lfloor 1, \pi_h(m) \rfloor \) is a \(\Lambda(A)\)-module homomorphism, and \((\Lambda(\pi), \Lambda(M))\), equipped with \( T_\Lambda \), is an object of   \(\textbf{Dil}_{\lambda}(A)\). Moreover, if 
    $\varphi : (\pi _1, M_1) \to  (\pi _2, M_2) $ is morphism of covariant representations of \(\lambda\), then the map $\Lambda (\varphi ) : \Lambda (M_1) \to \Lambda (M_1),$ such that $\Lambda (\varphi )(\lfloor g,m \rfloor )= \lfloor g, \varphi (m) \rfloor,$
    for any $\in \Lambda (M_1),$ is a morphism 
    \((\Lambda(\pi), \Lambda(M)) \to (\Lambda(\pi), \Lambda(M))\) in the category \(\textbf{Dil}_{\lambda}(A)\).  Consequently, \(\Lambda\) defines a functor from  \(\textbf{CovRep}_{\lambda} (A)\) into \(\textbf{Dil}_{\lambda}(A)\).
\end{lemma}

\begin{theorem}
    Let \(\lambda: G \to \operatorname{End}_0(A)\) be  an ideal  partial representation  of \(G\) on \(A\). Then the functor 
    \(
    \Lambda: \textbf{CovRep}_{\lambda}(A) \to \textbf{Dil}_{\lambda}(A) 
    \)
    is left adjoint to the functor 
    \(
    \F: \textbf{Dil}_{\lambda}(A) \to \textbf{CovRep}_{\lambda}(A).
    \)
\end{theorem}

\begin{proof}
    Let \((\pi, M)\) be an object in \(\textbf{CovRep}_{\lambda}(A)\), and let \((\rho, W)\), equipped with the \(\Lambda(\lambda)\)-module homomorphism \(T: W \to W\), be an object in \(\textbf{Dil}_{\lambda}(A)\). Recall, from Lemma~\ref{l: T for Lambda M}, that \((\Lambda(\pi), \Lambda(M))\) is equipped with the \(\Lambda(A)\)-module homomorphism \(T_\Lambda: \Lambda(M) \to \Lambda(M)\) defined by \(\lfloor g, m \rfloor \mapsto \lfloor 1, \pi_g(m) \rfloor\).

    Define the mapping
    \[
        \eta: \Hom((\Lambda(\pi), \Lambda(M)), (\rho, W)) \to \Hom((\pi, M), (\xi, T(W))),
    \]
    which sends a morphism \(\gamma: \Lambda(M) \to W\) to the morphism \(\eta(\gamma): M \to T(W)\) given by \(m \mapsto \gamma(\lfloor 1, m \rfloor)\). We first verify that \(\eta\) is well-defined, i.e., that \(\eta(\gamma): M \to T(W)\) is a morphism of  dilated covariant representations of \(\lambda\). 
   Note that
    $$T \eta(\gamma)(m) = T \gamma (\lfloor 1, m \rfloor) = \gamma T_{\Lambda}  (\lfloor 1, m \rfloor) = \gamma   (\lfloor 1, m \rfloor) =  \eta(\gamma)(m) $$ for any $m\in M$  so that
    $ \eta(\gamma)(M) \subseteq T(W).$ 
    Next, for  \(x \in A\), \(m \in M\), and \(g \in G\), we compute
    \[
        \eta(\gamma)(xm) = \gamma(\lfloor 1, xm \rfloor) = \gamma(\lfloor 1, x \rfloor \lfloor 1, m \rfloor) = \lfloor 1, x \rfloor \gamma(\lfloor 1, m \rfloor) = x\gamma(\lfloor 1, m \rfloor)
        = x (\eta(\gamma)(m)),
    \]
    showing that \(\eta(\gamma)\) is a morphism of \(A\)-modules. Now, for \(g \in G\),
\begin{multline*}
\eta(\gamma)(\pi_g(m)) = \gamma(\lfloor 1, \pi_g(m) \rfloor) = \gamma(T_\Lambda(\lfloor g, m \rfloor)) =\\= T\gamma(\Lambda(\lambda)_g(\lfloor 1, m \rfloor))
= T\rho_g\gamma(\lfloor 1, m \rfloor) = \xi_g(\eta(\gamma)(m)),
\end{multline*}    
    proving that \(\eta(\gamma)\) is a morphism of partial representations  of $G.$

    Next, define the mapping
    \[
        \delta: \Hom((\pi, M), (\xi, T(W))) \to \Hom((\Lambda(\pi), \Lambda(M)), (\rho, W)),
    \]
    sending a morphism \(\psi: M \to T(W)\) to the morphism \(\delta(\psi): \Lambda(M) \to W\) defined by \(\lfloor h, m \rfloor \mapsto \rho_h\psi(m)\). We verify that \(\delta\) is well-defined, i.e., that \(\delta(\psi)\) is a morphism of covariant representations of \(\Lambda(\lambda)\). For \(\lfloor g, x \rfloor \in \Lambda(A)\) and \(\lfloor h, m \rfloor \in \Lambda(M)\), we have
    \begin{multline*}
        \delta(\psi)(\lfloor g,x\rfloor \lfloor h,m\rfloor)=\delta(\psi)(\lfloor g,x\pi_{g^{-1}h}(m)\rfloor)=\rho_g \psi(x\pi_{g^{-1}h}(m))=\rho_g(x\xi_{g^{-1}h}\psi(m))= \\ =\rho_g(xT\rho_{g^{-1}h}\psi(m)) \overset{\scriptsize\eqref{eq xw}}{=} \rho_g(\lfloor 1,x\rfloor\rho_{g^{-1}h}\psi(m))=\Lambda(\lambda)_g(\lfloor 1,x\rfloor) \rho_h\psi(m)=\lfloor g,x\rfloor \delta(\psi)(\lfloor h,m\rfloor),
    \end{multline*}
    proving that \(\delta(\psi)\) is a morphism of \(\Lambda(A)\)-modules. Additionally,
    \[
        \delta(\psi)(\Lambda(\pi)_g(\lfloor h, m \rfloor)) = \delta(\psi)(\lfloor gh, m \rfloor) = \rho_{gh}\psi(m) = \rho_g(\rho_h\psi(m)) = \rho_g\delta(\psi)(\lfloor h, m \rfloor),
    \]
    verifying that \(\delta(\psi)\) is a morphism of covariant representations. Finally, we check that \(\delta(\psi) \circ T_\Lambda = T \circ \delta(\psi)\):
\begin{multline*}
\delta(\psi)(T_\Lambda(\lfloor h, m \rfloor)) = \delta(\psi)(\lfloor 1, \pi_h(m) \rfloor) = \psi(\pi_h(m)) =\\= \xi_h(\psi(m)) = T\rho_h(\psi(m)) = T\delta(\psi)(\lfloor h, m \rfloor).
\end{multline*}    
A straightforward computation shows that \(\delta\) is the inverse of \(\eta\), and thus \(\eta\) is a bijection.

    To conclude the proof of the theorem, we must show that the bijection \(\eta\) is natural. This requires verifying that for any morphism \(\gamma: M \to N\) in  \(\textbf{CovRep}_{\lambda}(A)\)  and any morphism \(\psi: W \to V\) in  \(\textbf{Dil}_{\lambda}(A)\), the  diagrams 
    \[
        \begin{tikzcd}
            \Hom((\Lambda({\pi}'), \Lambda(N)), (\rho, W)) 
            \arrow[r, "\Lambda(\gamma)^*"] 
            \arrow[d, "\eta"'] 
            & \Hom((\Lambda(\pi), \Lambda(M)), (\rho, W))  
            \arrow[d, "\eta"']  \\
            \Hom(({\pi}', N), (\xi, T(W))) 
            \arrow[r, "\gamma^*"] 
            & \Hom((\pi, M), (\xi, T(W)))  
        \end{tikzcd}
    \]
    
    
    and 
    
   \[
        \begin{tikzcd}
           \Hom((\Lambda (\pi), \Lambda(M)), (\rho, W)) 
            \arrow[r, "\psi_*"] 
            \arrow[d, "\eta"'] 
            & \Hom((\Lambda (\pi), \Lambda(M)), (\rho', V)) 
            \arrow[d, "\eta"] \\
            \Hom((\pi , M), (\xi, T(W))) 
            \arrow[r, "\F(\psi)_*"] 
            & \Hom((\pi , M), (\xi', T(V)))
        \end{tikzcd}
    \]  
     are commutative, where 
     \begin{align*}
     &\Lambda(\gamma)^* = \Hom (\Lambda (\gamma ), (\rho , W)), 
     {\gamma} ^*= \Hom (\gamma , (\xi , T(W))),\\
     &\psi _* = \Hom ((\Lambda(\pi), \Lambda (M)), \psi), \F(\psi)_* = \Hom ((\pi,M), \F (\psi) ) .
    \end{align*} 
    
%In these diagrams, the upper stars denote %post-composition, while the lower stars represent %pre-composition.

   For let \(\phi \in \Hom((\Lambda({\pi}'), \Lambda(N)), (\rho, W))\). For any \(m \in M\), we have
    \[
        \eta(\Lambda(\gamma)^*(\phi))(m) = \eta(\phi \circ \Lambda(\gamma))(m) = (\phi \circ \Lambda(\gamma))(\lfloor 1, m \rfloor) = \phi(\lfloor 1, \gamma(m) \rfloor),
    \]
    and
    \[
        \gamma^*(\eta(\phi))(m) = (\eta(\phi) \circ \gamma)(m) = \eta(\phi)(\gamma(m)) = \phi(\lfloor 1, \gamma(m) \rfloor),
    \]
     showing  the commutativity of the  first diagram. Now, let \(\varphi \in \Hom((\Lambda(\pi), \Lambda(M)), (\rho, W))\). For any \(m \in M\), we see that
    \[
        \eta(\psi_*(\varphi))(m) = \eta(\psi \circ \varphi)(m) = (\psi \circ \varphi)(\lfloor 1, m \rfloor) = \psi(\varphi(\lfloor 1, m \rfloor)),
    \]
    and
    \[
        \F(\psi)_*(\eta(\varphi))(m) = (\F(\psi) \circ \eta(\varphi))(m) = (\psi|_{T(W)} \circ \eta(\varphi))(m) = \psi(\varphi(\lfloor 1, m \rfloor)),
    \]
    proving that the second diagram is also commutative.
\end{proof}

\end{document}
